% Complete Hadamard / Mittag-Leffler expansion of I_B in d=2.
% The expansion is exact. A universal multi-digit coefficient is not produced.
\documentclass[11pt]{article}
\usepackage[margin=1.05in]{geometry}
\usepackage{amsmath,amssymb,amsthm}
\usepackage{enumitem,booktabs}
\usepackage[colorlinks=true,linkcolor=blue,citecolor=blue,urlcolor=blue]{hyperref}

\newtheorem{proposition}{Proposition}[section]
\newtheorem{theorem}{Theorem}[section]
\theoremstyle{definition}
\newtheorem*{admission}{Admission}

\newcommand{\pO}{\textbf{[O]}}
\newcommand{\pA}{\textbf{[ADMISSION]}}
\newcommand{\IB}{I_B}
\newcommand{\sech}{\operatorname{sech}}

\title{\textbf{Completion of the Hadamard Expansion of $\IB$\\
in Two Dimensions}}
\author{Robert W.\ McGwier\\
\small CTO, Cohere Technology Group\\
\small GitHub: \texttt{n4hy}}
\date{15 August 2026 --- Version 1}

\begin{document}
\maketitle

\begin{abstract}
\noindent
The multi-digit coefficient of the kinematic residue $\IB$ was left
open pending an analytic Hadamard expansion of the current two-point
function under modular flow. That expansion is completed here in
$d=2$.

The identity
\begin{equation*}
\frac{1}{\sinh^{2}w}
=
\sum_{n\in\mathbb{Z}}\frac{1}{(w-in\pi)^{2}}
\end{equation*}
is exact off the poles $w\in i\pi\mathbb{Z}$. The $n=0$ term is the
coincidence singularity. The Hadamard-finite remainder is the $n\neq 0$
sum, equivalently $1/\sinh^{2}w-1/w^{2}$. Folding against the modular
kernel reduces to a polygamma kernel $H(\tau)$ derived below.
The formula for the building block $I_{\mathrm{sing}}(\zeta)$ matches
an independent $\sech^{2}$ quadrature to $2\times 10^{-13}$.

$H(\tau)$ is \emph{not} proportional to $1/\tau^{2}$, to
$1/\sinh^{2}\tau$, to $I_{\mathrm{sing}}(\tau)$, or to
$1/\sinh^{2}\tau-1/\tau^{2}$. Relative residues formed with four
Gaussian sources therefore move with the source. There is no single
universal multi-digit coefficient of $\IB$ independent of $\beta$.

The Final Status label of $\IB$ stays \pO{} if ``coefficient'' means
one number. The expansion itself is closed.
\end{abstract}

\section*{Rule}
This note distinguishes an exact identity, a machine-checked special
function formula, a numerical spread, and an admission. It does not
upgrade a Final Status \pO{} to a derived number that the algebra
does not contain.

\section{Setup}

On the interval $(-R,R)$ with $R=1$, modular flow is
$x(s)=\tanh(\mathrm{artanh}\,x+\pi s)$. After $x=\tanh u$ the
weight-one Jacobians cancel \cite{McGwierVI,McGwier12} and
\begin{equation}
Q(s)
=
\int d\tau\,C(\tau)\,\frac{1}{\sinh^{2}(\tau+\pi s)},
\qquad
C(\tau)=\int du\,b(u)b(u-\tau),
\quad
b(u)=\beta(\tanh u).
\label{eq:Q}
\end{equation}
The kernel is $K(s)=\pi/(2\cosh^{2}(\pi s))$, $\int K=1$, and
\begin{equation}
\int K(s)\,f(\tau+\pi s)\,ds
=
\int\frac{d\alpha}{2\cosh^{2}\alpha}\,f(\tau+\alpha).
\label{eq:pull}
\end{equation}

\section{The complete expansion}

\begin{theorem}[Mittag--Leffler]
For $w\notin i\pi\mathbb{Z}$,
\begin{equation}
\frac{1}{\sinh^{2}w}
=
\sum_{n=-\infty}^{\infty}\frac{1}{(w-in\pi)^{2}}.
\label{eq:ML}
\end{equation}
\end{theorem}

This is the Hadamard series: every pole is written, none is omitted.
The coincidence is $n=0$. The finite remainder
\begin{equation}
h(w)
:=
\frac{1}{\sinh^{2}w}-\frac{1}{w^{2}}
=
\sum_{n\neq 0}\frac{1}{(w-in\pi)^{2}}
\label{eq:h}
\end{equation}
is entire in a neighbourhood of the real axis and equals $-1/3$ at
$w=0$.

The Hadamard-finite smeared correlator is therefore
\begin{equation}
Q_{H}(s)
=
\int d\tau\,C(\tau)\,h(\tau+\pi s)
=
\sum_{n\neq 0}\int d\tau\,\frac{C(\tau)}{(\tau+\pi s-in\pi)^{2}}.
\label{eq:QH}
\end{equation}

\section{The modular integral of the finite part}

Define
\begin{equation}
I_{\mathrm{sing}}(\zeta)
:=
\frac12\int_{-\infty}^{\infty}d\alpha\,
\frac{\sech^{2}\alpha}{(\alpha+\zeta)^{2}}.
\label{eq:Ising}
\end{equation}
Residues of $\sech^{2}$ at $i\pi(n+\tfrac12)$ give, for
$\operatorname{Im}\zeta>0$,
\begin{equation}
I_{\mathrm{sing}}(\zeta)
=
\pi^{-2}\,\psi^{(2)}\!\left(\tfrac12-\frac{i\zeta}{\pi}\right),
\label{eq:pg}
\end{equation}
and $I_{\mathrm{sing}}(-\zeta)=I_{\mathrm{sing}}(\zeta)$. On the real
axis
$I_{\mathrm{sing}}(\tau)=\pi^{-2}\operatorname{Re}\psi^{(2)}(\tfrac12-i\tau/\pi)$.

Machine check (\texttt{m9\_4\_ib\_hadamard\_complete.py}):
\eqref{eq:pg} versus $\sech^{2}$ quadrature at three complex points,
maximum relative error $1.99\times 10^{-13}$.

The finite modular integral is the kernel
\begin{align}
H(\tau)
&:=
\int K(s)\,h(\tau+\pi s)\,ds
=
\sum_{n\neq 0}I_{\mathrm{sing}}(\tau-in\pi)
\nonumber\\
&=
\pi^{-2}
\sum_{m=1}^{\infty}
\Bigl[
\psi^{(2)}\!\bigl(m+\tfrac12+\tfrac{i\tau}{\pi}\bigr)
+
\psi^{(2)}\!\bigl(m+\tfrac12-\tfrac{i\tau}{\pi}\bigr)
\Bigr].
\label{eq:H}
\end{align}
Then
\begin{equation}
\int K(s)\,Q_{H}(s)\,ds
=
\int d\tau\,C(\tau)\,H(\tau).
\label{eq:IQ}
\end{equation}

\section{There is no universal ratio}

A single coefficient would mean $H=\lambda K_{\mathrm{loc}}$ for some
local kernel $K_{\mathrm{loc}}$ independent of $C$. On
$\tau\in\{0.3,0.5,0.8,1.0,1.5,2.0,3.0\}$ the ratio $H/K_{\mathrm{loc}}$
has spread

\begin{center}
\begin{tabular}{@{}lc@{}}
\toprule
$K_{\mathrm{loc}}$ & spread of $H/K_{\mathrm{loc}}$ \\
\midrule
$I_{\mathrm{sing}}(\tau)$ & $0.553$ \\
$1/\tau^{2}$ & $0.968$ \\
$h(\tau)$ & $0.210$ \\
$1/\sinh^{2}\tau$ & not constant \\
\bottomrule
\end{tabular}
\end{center}

Relative residues $r=\int CH/\int C K_{\mathrm{loc}}$ for Gaussians
of width $\sigma\in\{0.4,0.7,1.0,1.4\}$:

\begin{center}
\begin{tabular}{@{}lcc@{}}
\toprule
denominator & mean $r$ & spread in $\sigma$ \\
\midrule
$I_{\mathrm{sing}}$ & $0.0585$ & $7.08\times 10^{-2}$ \\
$h$ & $0.898$ & $4.49\times 10^{-2}$ \\
$1/\tau^{2}$ & $-0.0203$ & $1.64\times 10^{-2}$ \\
$1/\sinh^{2}$ & $-0.0208$ & $1.71\times 10^{-2}$ \\
\bottomrule
\end{tabular}
\end{center}

None of these spreads is a rounding error.

\begin{admission}
The Hadamard series is complete. What it does \emph{not} produce is
one multi-digit number independent of $\beta$. Asking for that number
was slightly the wrong question: the finite object is the kernel
$H(\tau)$, and $H$ is not a multiple of a local quadratic form.
The Final Status line ``multi-digit coefficient of $\IB$'' remains
\pO{} as a \emph{number}. As an expansion, the $d=2$ problem is
closed.
\end{admission}

\section{What remains}

\begin{itemize}[leftmargin=1.4em]
\item $d=4$ balls, where the current two-point function is not
$1/(x-y)^{2}$ on a line and CHM is a conformal image of Rindler, not
a M\"obius map on $\mathbb{R}$.
\item A renormalization convention that would force a unique
$\lambda$ by extra physical input. None is supplied by $K(s)$ alone.
\item Digits of $H(\tau)$ at a conventional $\tau$ are available from
\eqref{eq:H} to any precision; they are values of a function, not a
coefficient of $\IB$.
\end{itemize}

\section*{Script}
\texttt{m9\_4\_ib\_hadamard\_complete.py}, data
\texttt{m9\_4\_ib\_hadamard\_complete.json}.

\begin{thebibliography}{99}
\bibitem{McGwierVI} R.W.~McGwier, ``Evaluation of the Universal Kinematic Integral for Condition NL,''
\url{https://github.com/n4hy/New_Model_Emergent_Gravity}.
\bibitem{McGwierVIII} R.W.~McGwier, ``Status of the High-Precision Coefficient of $\IB$ (v4),'' ibid.
\bibitem{McGwier12} R.W.~McGwier, ``Analytic Attempt at the Coefficient of $\IB$,'' ibid.
\bibitem{CHM} H.~Casini, M.~Huerta and R.C.~Myers, JHEP 05 (2011) 036.
\bibitem{CTT} H.~Casini, E.~Teste and G.~Torroba, J.\ Phys.\ A 50 (2017) 364001.
\end{thebibliography}
\end{document}
